\section{The original neat levels for every auxiliary level}
\label{sec:neat-levels}

This chapter proves the full predicate \texttt{UpperTheory.V.Hi.neatU}.
Its conclusion concerns the actual frame and level subgroups used to
construct ordinary automorphic forms.  The integral representation
\(\tau _0\) is arbitrary.  Neither a type theorem nor an ordinary Hecke
axiom is an input to the proof.

Here and below the source is pinned to commit
\texttt{7cbf7a0c32fca1badf05931879a37c3f8a2b03b4}.
For source references put
\[
 \mathsf H=\texttt{HeckeAlgebras/},\qquad
 \mathsf O=\texttt{HeckeAlgebras/OrdinaryHeckeCharacters/},
\]
under
\texttt{src/Development/AutomorphicArithmetic/AutomorphicVarieties/}.
The main source is \(\mathsf O\)\texttt{NtExpand.lean}; precise companion
definitions and inclusive source ranges are recorded with every claim
in the accompanying provenance file.  The Lean proof is accepted as
correct.  The arguments below establish the human mathematical endpoint
independently of invoking that proof.

\subsection{The datum and the exact conclusion}

Let \(\ell\) be a prime, let \(L\) be a CM number field with real subfield
\(F=L^+\), and write \(c_L\) for its complex conjugation.  Fix the original
algebraic closure \(\Omega\) and the original
\(s:\texttt{HidaSetting}\ \ell\ L\ \Omega\).  Among its fields are a
choice \(w(v)\) of a place of \(L\) above each finite place \(v\) of \(F\),
a finite set \(S_t\) of places away from \(\ell\), and a distinguished
place \(v_0\notin S_t\).  Its auxiliary rational prime \(p=s.p\) satisfies
\begin{equation}\label{eq:ntl-aux-data}
 p\text{ is prime},\quad p\ne2,\quad p\ne\ell,\quad
 p\nmid\operatorname{disc}(L),\quad
 \#(\mathcal O_F/v_0)=p.
\end{equation}
This is the prime already stored in \(s\).  It is not a prime selected
as a function of the matrix rank.

Put
\(S=\{v:v\mid\ell\}\cup S_t\cup\{v_0\}\).
Fix \(n\ge1\), an original \(V:s.\texttt{Lev}\), and
\(Q:\texttt{QData}\ V\).  Thus \(Q\) has a finite index set, places
\(w_q\) of \(L\), and distinct underlying places
\(v_q=w_q\cap\mathcal O_F\), with
\begin{equation}\label{eq:ntl-Q}
 v_q\notin S,\qquad V.S'=S\cup\{v_q:q\in Q\}.
\end{equation}
Write \(G_w=\operatorname{GL}_n(L_w)\),
\(K_w^{\max}=\operatorname{GL}_n(\mathcal O_{L_w})\), viewed in \(G_w\),
and
\[
 C_w(1)=\ker\bigl(\operatorname{GL}_n(\mathcal O_{L_w})
       \longrightarrow\operatorname{GL}_n(k_w)\bigr),
\]
again viewed in \(G_w\).  These are the source \texttt{KO} and
\texttt{Kc}\ \(1\), respectively.

The source group \(G^a\) is the image of the adelic definite unitary
group under all its projections at places of \(L\) away from \(\ell\).
For every place \(w\), its homomorphism \(\operatorname{cmp}_w\) is that
coordinate when \(w\nmid\ell\), and is the trivial map otherwise.
Its subgroup \(U_b\) imposes membership in \(K_{w(u)}^{\max}\) for
\(u\in S_t\), and in \(C_{w(v_0)}(1)\) at the distinguished place.
The coefficient representation is any homomorphism
\[
 \tau _0:\prod_{u\in S_t}K_{w(u)}^{\max}
                  \longrightarrow\operatorname{GL}_m(\mathcal O),
 \qquad \mathcal O=c.\mathcal O
\]
for the original \(c:\texttt{Coef}\ \ell\) and any \(m\ge0\).
Its value does not enter the group argument.

Let \(H=\texttt{hset}\ Q\ \tau _0\), \(X=H.X\), and
\(\Gamma=\texttt{ratU}\ L\ n\).  The frame has
\(G^\ell=\prod_{v\mid\ell}G_{w(v)}\) and the original rational map
\(\operatorname{glob}:\Gamma\longrightarrow G^a\times G^\ell\).
For each family of compact open subgroups \(K=(K_q)_{q\in Q}\),
the actual level is
\begin{equation}\label{eq:ntl-level}
 U(K)=U_Q\vee\bigvee_{q\in Q}\operatorname{sec}_{w_q}(K_q)
                  \quad\text{inside }U_b.
\end{equation}
Here \(U_Q\) is integral outside \(V.S'\) and has trivial \(w_q\)
coordinate for every \(q\).  The sections in this formula are exactly
the source sections, including their trivial-map fallback.

The desired \texttt{Neat} assertion is the following implication:
\begin{equation}\label{eq:ntl-neat}
 \begin{gathered}
 g\in G^a\times G^\ell,\quad\gamma\in\Gamma,\quad
 u\in U(K),\quad b\in X.I.\operatorname{Iw}(0),\\
 \operatorname{glob}(\gamma)g=g(u,b)
       \quad\Longrightarrow\quad u=1\text{ and }b=1.
 \end{gathered}
\end{equation}
Every object in this formula is the original one.  In particular,
\texttt{Neat} is used in its displayed source sense; no stronger
definition involving all eigenvalue-generated subgroups is substituted.

\subsection{Torsion in the actual principal congruence subgroup}

% BEGIN NTL1
\paragraph{NTL1. A third-order identity in a possibly noncommutative ring.}
\label{claim:NTL1}
For every unital ring \(R\), every \(Y\in R\), and every integer
\(a\ge0\), there is \(W\in R\) such that
\begin{equation}\label{eq:ntl-expand}
 (1+Y)^a=1+aY+\binom a2Y^2+Y^3W.
\end{equation}
Only powers of the single element \(Y\) occur, so their usual binomial
expansion is valid without commutativity of \(R\).  Alternatively,
start with \(W=0\) for \(a=0\).  Multiplication by \(1+Y\) changes the
remainder to \(W(1+Y)+\binom a2\), and
\(\binom{a+1}2=a+\binom a2\) gives the formula at the next exponent.
This proves the identity also for matrix rings of every size.
\emph{Source:} \texttt{nt\_expand}, lines 12--24.
% END NTL1

% BEGIN NTL2
\paragraph{NTL2. One extra power of the auxiliary prime.}
\label{claim:NTL2}
Let \(R\) be a commutative local domain, and let \(p\) be an odd
rational prime whose image in \(R\) is nonzero and is not a unit.
For any rational prime \(q\), any \(k\ge0\), and any matrix \(Z\),
\[
 (1+p^{k+1}Z)^q=1\quad\Longrightarrow\quad
                    Z\in p\operatorname{Mat}_n(R).
\]
Apply NTL1 and cancel the nonzero central scalar \(p^{k+1}\)
entry by entry.  If \(q\ne p\), Bezout's identity for \(q,p\) shows
that \(q\) is a unit in the local ring: \(p\) belongs to the maximal
ideal, so \(q\) cannot belong to it.  The resulting identity is
\[
 qZ=-p\left(\binom q2p^kZ^2+p^{2k+1}Z^3W\right),
\]
which proves the assertion.  If \(q=p\), write
\(\binom p2=pc\) with \(c=(p-1)/2\in\mathbb Z\).
Cancel \(p^{k+2}\) in the original expansion to obtain
\[
 Z=-p\left(cp^kZ^2+p^{2k}Z^3W\right).
\]
This proves the second case.  No coefficient depending on the rank,
and no factor \(n!\), is inverted.
\emph{Source:} \texttt{nt\_step}, lines 28--60.
% END NTL2

% BEGIN NTL3
\paragraph{NTL3. Torsion-freeness over a Noetherian local domain.}
\label{claim:NTL3}
Under the hypotheses of NTL2, assume in addition that \(R\) is
Noetherian.  Every finite-order matrix in
\(1+p\operatorname{Mat}_n(R)\) is the identity.

First, \(\bigcap_{a\ge0}p^aR=0\).  To see this directly, if a nonzero
\(x\) belongs to every \(p^aR\), put \(z_a=x/p^a\in R\), using the
fraction field.  The ideals \((z_a)\) form an increasing chain because
\(z_a=pz_{a+1}\).  At a stationary pair,
\(z_{a+1}=r z_a\) for some \(r\in R\), whence
\(z_a=prz_a\).  Cancellation gives \(pr=1\), a contradiction.

Suppose now that \(A=1+pZ\) satisfies \(A^q=1\) for a prime \(q\).
Repeated application of NTL2 writes \(A-1\) in
\(p^{k+1}\operatorname{Mat}_n(R)\) for every \(k\ge0\).
The preceding separation argument makes all its entries zero.
For a general positive exponent \(a\) with \(A^a=1\), induct on \(a\).
If \(a=1\), there is nothing to prove.  Otherwise choose a prime
\(q\mid a\) and put \(a=qd\).  The matrix \(A^d\) is still congruent
to \(1\) modulo \(p\); its \(q\)-th power is \(1\), so the prime case
gives \(A^d=1\).  Since \(d<a\), induction finishes the proof.
\emph{Source:} \texttt{nt\_prime} and \texttt{nt\_mat}, lines 62--99.
% END NTL3

% BEGIN NTL4
\paragraph{NTL4. Why the stored discriminant condition makes \(p\) a uniformizer.}
\label{claim:NTL4}
Let \(E\) be a number field, let \(p\) be a rational prime with
\(p\nmid\operatorname{disc}(E)\), and let \(w\mid p\).  Then
\[
 p\mathcal O_{E_w}=\mathfrak m_{E_w},
 \qquad p\notin w^2\subset\mathcal O_E.
\]
Here is a trace-form proof of the discriminant step.  Choose an integral
basis of \(\mathcal O_E\).  Its trace-pairing matrix has determinant
\(\operatorname{disc}(E)\), so its reduction modulo \(p\) is
nondegenerate on the finite-dimensional commutative algebra
\(A=\mathcal O_E/p\mathcal O_E\).  This is the trace pairing for the
regular representation of \(A\): multiplication matrices reduce
coefficientwise modulo \(p\).  If \(x\in A\) is nilpotent, then
\(xy\) is nilpotent for every \(y\in A\), and the multiplication
operator by \(xy\) is nilpotent.  Its trace is zero.  Nondegeneracy
therefore forces \(x=0\).

A finite reduced commutative algebra over a finite field is a product
of fields.  Indeed it is Artinian; its finitely many maximal ideals
are pairwise comaximal, and their intersection is its nilradical,
which is zero.  The Chinese remainder map is consequently an
isomorphism onto the product of its residue fields.  Localizing at
the maximal ideal corresponding to \(w\) now shows that
\(\mathcal O_{E,w}/p\mathcal O_{E,w}\) is a field.  Hence
\(p\mathcal O_{E,w}\) is its maximal ideal.  The local ring is a
discrete valuation ring, so the valuation of \(p\) is exactly one.
Completion preserves this equality of the maximal ideal with
\((p)\), and membership of a global element in \(w^2\) would give
local valuation at least two, proving the assertions.

The elementary number-field structure used here does not conceal a
prime-selection theorem.  In particular, finiteness of
\(\mathcal O_E\) as a free \(\mathbb Z\)-module can be seen by taking
an integral \(\mathbb Q\)-basis \(e_i\) of \(E\).  For
\(x\in\mathcal O_E\), the numbers \(\operatorname{Tr}(xe_i)\) are
integers; the invertible rational trace matrix therefore puts every
such \(x\) in the dual of one fixed lattice.  This bounds denominators
and gives a finite free integral basis.

For completeness, the local discrete-valuation assertion also has
an elementary algebraic proof.  The ring of integers is Noetherian
by this finiteness, and is integrally closed by transitivity of
integrality.  Every nonzero prime contracts to a nonzero prime of
\(\mathbb Z\): a nonzero element in it has a monic integer equation
with nonzero constant coefficient, which puts that coefficient in
the contraction.  Its quotient is a finite integral domain, hence
a field.  Thus a nonzero-prime localization \(R\) is a Noetherian
normal local domain of dimension one.  Choose \(0\ne a\in\mathfrak m\).
There is \(r\ge1\) with \(\mathfrak m^r\subseteq(a)\), since
\(\sqrt{(a)}=\mathfrak m\) and \(\mathfrak m\) is finitely generated.
Take the least such \(r\).  If \(r=1\), then \(\mathfrak m=(a)\).
Otherwise choose \(b\in\mathfrak m^{r-1}\setminus(a)\), and put
\(t=b/a\).  Then \(t\notin R\) and \(t\mathfrak m\subseteq R\).
If \(t\mathfrak m\subseteq\mathfrak m\), express multiplication
by \(t\) on a finite generating list of \(\mathfrak m\).
The adjugate identity gives a monic polynomial \(P\in R[T]\) with
\(P(t)\mathfrak m=0\), hence \(P(t)=0\) in the fraction field
because \(\mathfrak m\ne0\).  Normality would
give \(t\in R\), a contradiction.  Therefore \(t\mathfrak m\)
contains a unit.  After rescaling its preimage, there is
\(x\in\mathfrak m\) with \(tx=1\), and every
\(y\in\mathfrak m\) is \(x(ty)\).  Thus \(\mathfrak m=(x)\).
The increasing-principal-ideal argument in NTL3, with \(x\) in
place of \(p\), shows \(\bigcap_jx^jR=0\); every nonzero element
is consequently a unit times a finite power of \(x\).  This is the
discrete valuation description.  All uses here are at the already
supplied \(w\mid p\).
\emph{Source:} \texttt{nt\_sq}, lines 115--136, and the valuation
translation in \texttt{nt\_free}, lines 140--176.
% END NTL4

% BEGIN NTL5
\paragraph{NTL5. The exact local \texttt{Kc}\ \(1\) has no nontrivial torsion.}
\label{claim:NTL5}
Let \(E,p,w\) satisfy NTL4, assume \(p\ne2\), and let
\(g\in G_w\) lie in the original \(C_w(1)\).
If \(g\) has finite order, then \(g=1\).

The definition of \(C_w(1)\) gives a matrix
\(x\in\operatorname{GL}_n(\mathcal O_{E_w})\) mapping to \(g\) and
reducing to the identity modulo \(\mathfrak m_{E_w}\).
By NTL4, this is exactly the condition \(x=1+pZ\) for a matrix
\(Z\) over \(\mathcal O_{E_w}\).  The inclusion of the integral
matrix group into \(G_w\) is injective, so a positive exponent killing
\(g\) also kills \(x\).  The ring \(\mathcal O_{E_w}\) is a
Noetherian local domain, its element \(p\) is nonzero, and \(p\) is
not a unit.  NTL3 applies and gives \(x=1\), hence \(g=1\).
This treats precisely reduction modulo the maximal ideal; it would
not be justified by replacing that ideal with \((p)\) without NTL4.
\emph{Source:} \texttt{nt\_free}, lines 138--203; the definition of
\texttt{Kc} is \(\mathsf H\)\texttt{TcAnti.lean}, lines 69--71.
% END NTL5

\subsection{From local compactness to finite order over the CM field}

% BEGIN NTL6
\paragraph{NTL6. A positive power of any compact-level element is integral.}
\label{claim:NTL6}
If \(K\) is a compact subgroup of \(\operatorname{GL}_n(E_w)\) and
\(g\in K\), then there is an integer \(e>0\) such that
\(g^e\in K_w^{\max}\).

The subgroup \(K_w^{\max}\) is open: a sufficiently small perturbation
of an integral invertible matrix is again integral with unit
determinant, as is seen by reduction modulo the maximal ideal.
Consequently \(K\cap K_w^{\max}\) is open in the compact group \(K\).
Its left cosets form an open cover, so only finitely many occur.
Two of the cosets represented by \(1,g,g^2,\ldots\) coincide.  If
their indices are \(a<b\), then \(g^{b-a}\in K\cap K_w^{\max}\).
This is a coset argument; no normality of \(K\cap K_w^{\max}\) is
needed.  The exponent may depend on this element and this compact
subgroup, exactly as in the source.
\emph{Source:} \texttt{nt\_pow}, lines 205--212.
% END NTL6

% BEGIN NTL7
\paragraph{NTL7. Integral local conjugates give integral characteristic coefficients.}
\label{claim:NTL7}
Let \(M\in\operatorname{GL}_n(E)\), and suppose that for some
\(g_w\in\operatorname{GL}_n(E_w)\),
\[
 g_w^{-1}M_wg_w\in K_w^{\max}.
\]
Every coefficient of \(\det(T-M)\) is integral at \(w\).  Indeed the
displayed conjugate is represented by a matrix over
\(\mathcal O_{E_w}\); its characteristic coefficients are polynomials
with integral coefficients in its entries.  Characteristic
polynomials are unchanged by conjugation and commute with the field
map \(E\longrightarrow E_w\).  Thus the image of every coefficient
belongs to \(\mathcal O_{E_w}\), which is precisely the asserted
nonnegative valuation in \(E\).  The same statement applies to
\(M^{-1}\) with the same conjugating matrix, because
\(K_w^{\max}\) is a group.
\emph{Source:} \texttt{nt\_char}, lines 214--231.
% END NTL7

% BEGIN NTL8
\paragraph{NTL8. One place above each real-field place, with conjugation, is enough.}
\label{claim:NTL8}
For each finite place \(v\) of \(F=L^+\), choose any place \(T(v)\)
of \(L\) above \(v\).  If \(x\in L\) and both \(x\) and
\(c_L(x)\) are integral at every chosen \(T(v)\), then \(x\) is
an algebraic integer.

First, if \(y\in F\) is integral at every \(T(v)\), then it is
integral at every \(v\).  The restriction of the normalized valuation
at \(T(v)\) multiplies the valuation at \(v\) by a positive integer:
write an element of \(F^\times\) as a power of a local uniformizer
times a local unit.  The unit stays a unit upstairs, and the
uniformizer has positive valuation because the prime upstairs
contracts to \(v\).  The sign of the valuation is therefore unchanged.
Integrality at every finite place says \(y\in\mathcal O_F\).
For clarity, the local-to-global assertion follows directly from
\(\mathcal O_F=\bigcap_v(\mathcal O_F)_v\): if \(y\notin\mathcal O_F\),
the proper denominator ideal
\(I=\{a\in\mathcal O_F:ay\in\mathcal O_F\}\) lies in a maximal
ideal \(v\); an expression \(y=b/s\) in \((\mathcal O_F)_v\)
would give \(s\in I\) and \(s\notin v\), a contradiction.

Now \(a=x+c_L(x)\) and \(b=xc_L(x)\) are fixed by \(c_L\), hence
belong to \(F\).  They are integral at each \(T(v)\), by the ring
property of local integers, so the preceding paragraph gives
\(a,b\in\mathcal O_F\).  The element \(x\) satisfies
\(T^2-aT+b=0\).  It is therefore integral over \(\mathcal O_F\)
and hence over \(\mathbb Z\).  Transitivity here can also be read
from the finite \(\mathbb Z\)-module
\(\mathcal O_F+\mathcal O_Fx\), stable under multiplication by \(x\):
the characteristic polynomial of this multiplication is monic
with integer coefficients and kills \(x\).
\emph{Source:} \texttt{nt\_int}, lines 239--261.
% END NTL8

% BEGIN NTL9
\paragraph{NTL9. Bounded algebraic integers in a fixed number field form a finite set.}
\label{claim:NTL9}
Fix a number field \(E\) of degree \(d\) and a real number \(C\ge0\).
There are only finitely many \(x\in\mathcal O_E\) such that
\(\lvert\varphi(x)\rvert\le C\) for every embedding
\(\varphi:E\hookrightarrow\mathbb C\).

The minimal polynomial of such an \(x\) has degree \(r\le d\) and
integer coefficients.  Every embedding of \(\mathbb Q(x)\) into
\(\mathbb C\) extends to one of \(E\), so all \(r\) roots have
absolute value at most \(C\).  Its coefficient at \(T^{r-j}\)
has absolute value at most \(\binom rj C^j\), by its elementary
symmetric expression in those roots.  For each of the finitely
many possible degrees, these bounds leave only finitely many integer
coefficient lists.  Each resulting nonzero polynomial has finitely
many roots.  Their finite union contains all the possible \(x\).
This proves the finiteness used below without a covolume or
geometry-of-numbers estimate.
\emph{Source match:} the bounded-integral-entry step in
\texttt{nt\_fin}, lines 270--299, where the source uses
\texttt{Embeddings.finite\_of\_norm\_le}.
% END NTL9

% BEGIN NTL10
\paragraph{NTL10. Integral characteristic coefficients and definite unitarity force finite order.}
\label{claim:NTL10}
Suppose \(M\in\operatorname{GL}_n(L)\) has integral characteristic
coefficients and
\begin{equation}\label{eq:ntl-unitary-powers}
 c_L(M^a)^{\mathsf t}M^a=1\qquad(a\ge0).
\end{equation}
Then \(M\) has finite order.

Cayley--Hamilton expresses every nonnegative power of \(M\) as an
\(\mathcal O_L\)-linear combination of
\(1,M,\ldots,M^{n-1}\).  Choose a positive integer \(D\) clearing
the denominators of the entries of these finitely many matrices.
Such an integer exists for each algebraic number: after clearing
the rational coefficients of a monic polynomial over \(\mathbb Q\),
a suitable integer multiple of its root satisfies a monic polynomial
over \(\mathbb Z\).  Thus all entries of \(DM^a\), for all \(a\),
are algebraic integers.

For an embedding \(\varphi:L\hookrightarrow\mathbb C\), the CM
identity \(\varphi\circ c_L=\overline{\varphi}\) turns
\eqref{eq:ntl-unitary-powers} into ordinary complex unitarity of
\(\varphi(M^a)\).  Each entry has absolute value at most one: the
sum of squared absolute values in its column is one.  Every conjugate
of each entry of \(DM^a\) is therefore bounded by \(D\).
NTL9 leaves only finitely many values for each entry.  Since the
matrix has finitely many entries, \(\{M^a:a\ge0\}\) is finite.
Choose \(a<b\) with \(M^a=M^b\) and cancel \(M^a\).  Then
\(M^{b-a}=1\), with \(b-a>0\), as required.
\emph{Source:} \texttt{nt\_fin}, lines 263--302.
% END NTL10

% BEGIN NTL11
\paragraph{NTL11. A compact conjugate at one place above every \(v\) suffices.}
\label{claim:NTL11}
Let \(M\in\operatorname{GL}_n(L)\) satisfy
\eqref{eq:ntl-unitary-powers}.  Suppose that, for each finite place
\(v\) of \(F\), there are a place \(T(v)\mid v\) and a matrix
\(g_v\in G_{T(v)}\) for which
\[
 g_v^{-1}M_{T(v)}g_v\in K_{T(v)}^{\max}.
\]
Then \(M\) has finite order.

Let \(a_i\) be a characteristic coefficient of \(M\).  NTL7 makes
\(a_i\) integral at every \(T(v)\).  The case \(a=1\) of
\eqref{eq:ntl-unitary-powers} gives
\(M^{-1}=c_L(M)^{\mathsf t}\).  The characteristic coefficients of
this matrix are \(c_L(a_i)\), because transpose preserves the
characteristic polynomial and field automorphisms act on its
coefficients.  The inverse of the displayed compact conjugate lies
in the same \(K_{T(v)}^{\max}\), so NTL7 also makes
\(c_L(a_i)\) integral at every \(T(v)\).  NTL8 gives
\(a_i\in\mathcal O_L\).  Apply NTL10.

The choice \(T(v)\) need not be the original \(w(v)\) at every place.
This matters at the added places, where the original datum gives
\(w_q\) rather than an equality \(w_q=w(v_q)\).
\emph{Source:} \texttt{nt\_ord}, lines 304--322.
% END NTL11

\subsection{Literal frame and level identities}

% BEGIN NTL12
\paragraph{NTL12. Extracting the rational matrix from the original \(\Gamma\).}
\label{claim:NTL12}
For every \(\gamma\in\Gamma=\texttt{ratU}\ L\ n\), there is a
matrix \(M\in\operatorname{GL}_n(L)\) whose rational adelic image
is the underlying element of \(\gamma\), such that
\eqref{eq:ntl-unitary-powers} holds and, for every finite place \(w\),
\begin{equation}\label{eq:ntl-rational-coordinate}
 \operatorname{cLoc}_w(\gamma)=M_w.
\end{equation}
Existence of \(M\) is the definition of \texttt{ratU}: it is the
intersection of the rational matrix image with the adelic unitary
subgroup for the identity Hermitian matrix.  The latter condition
is \(\operatorname{conjAd}(\gamma)^{\mathsf t}\gamma=1\).
Adelic conjugation on an element of \(L\) is the image of \(c_L\).
Since the diagonal map from \(L\) into its adeles is injective,
the adelic equality descends entrywise to
\(c_L(M)^{\mathsf t}M=1\).  Its powers satisfy the same identity:
\(c_L(M^a)^{\mathsf t}=(c_L(M)^{\mathsf t})^a=(M^{-1})^a\).
Projection of the rational adelic image to the completion gives
\eqref{eq:ntl-rational-coordinate} by the definitions of
\texttt{ratGL} and \texttt{cLoc}.
\emph{Source:} \texttt{nt\_rat}, lines 334--351; the underlying
\texttt{UA}, \texttt{cLoc}, and \texttt{ratU} are defined in
\texttt{BaseChangeDefiniteUnitary.lean}, lines 66--84.
% END NTL12

% BEGIN NTL13
\paragraph{NTL13. Both required level inclusions, including the section fallback.}
\label{claim:NTL13}
For the actual \(U(K)\) in \eqref{eq:ntl-level}, the following hold:
\begin{align}
 v(w)\notin V.S'&\quad\Longrightarrow\quad
           \operatorname{cmp}_w(U(K))\subseteq K_w^{\max},
                 \label{eq:ntl-level-out}\\
 q\in Q&\quad\Longrightarrow\quad
           \operatorname{cmp}_{w_q}(U(K))\subseteq K_q.
                 \label{eq:ntl-level-Q}
\end{align}
Here \(v(w)=w\cap\mathcal O_F\).  To verify the inclusions on the
literal generators, recall that \(\operatorname{locA}_w\) is the
intersection of the kernels of all \(\operatorname{cmp}_{w'}\)
with \(v(w')\ne v(w)\).  By definition,
\(\operatorname{sec}_w\) is the inverse of the own-place projection
on this subgroup when that projection is bijective, and the trivial
homomorphism otherwise.  Hence for every \(x\in G_w\),
\begin{equation}\label{eq:ntl-sec}
 \begin{gathered}
 \operatorname{sec}_w(x)\in\operatorname{locA}_w,\\
 \operatorname{cmp}_w(\operatorname{sec}_w(x))=x
        \quad\hbox{or}\quad
 \operatorname{cmp}_w(\operatorname{sec}_w(x))=1.
 \end{gathered}
\end{equation}
If \(v(w)\notin V.S'\), the \(U_Q\) generator is integral at \(w\)
by definition; each \(\operatorname{sec}_{w_q}(K_q)\) has trivial
\(w\)-coordinate because \(v_q\in V.S'\).  This proves
\eqref{eq:ntl-level-out} on every generator.  For
\eqref{eq:ntl-level-Q}, the \(U_Q\) coordinate is trivial.
The generator at the same \(q\) has coordinate in \(K_q\) by
\eqref{eq:ntl-sec}, in both alternatives.  A generator at \(q'\ne q\)
has trivial coordinate, because the underlying places are distinct
by the original \(Q.\texttt{inj}\) condition.  Since the targets
\(K_w^{\max}\) and \(K_q\) are subgroups, both inclusions persist
under the subgroup joins.  This proves them on all of \(U(K)\).
No assertion of bijectivity for the source section was used.
\emph{Source:} \texttt{nt\_sec}, \texttt{nt\_out}, and
\texttt{nt\_Q}, lines 363--399; \(\mathsf H\)\texttt{IsTheta.lean},
lines 84--102 and 117--139.
% END NTL13

% BEGIN NTL14
\paragraph{NTL14. The original Iwahori subgroup is integral in every coefficient-prime coordinate.}
\label{claim:NTL14}
For the actual product Iwahori datum of \(H\),
\[
 X.I.\operatorname{Iw}(0)
             \ \subseteq\ \prod_{v\mid\ell}K_{w(v)}^{\max}.
\]
The source defines \(\operatorname{Iw}(0)\) as the subgroup generated
by \(N,T,N_m(0)\).  In each local factor, the defining formulas give
\(N=\operatorname{unipB}\cap K^{\max}\), so \(N\) is integral.
The subgroup \(T\) consists of diagonal matrices whose entries are
units of the local integer ring; the inverse diagonal matrix is
also integral.  Finally, \(N_m(0)\) is an intersection with
\(C_w(1)\), which lies in \(K^{\max}\) by its defining image from
the integral matrix group.  Thus all three local subgroups lie in
\(K^{\max}\).  The product definitions of \texttt{piData} and
\texttt{piIwa} give the same assertion coordinatewise, and the
product of the \(K^{\max}\)'s is a subgroup.  Taking the join proves
the assertion.  Full Iwahori factorization and the ordinary Hecke
operator identities are unnecessary for this inclusion.
\emph{Source:} \texttt{nt\_Iw}, lines 403--410; literal subgroup
fields in \(\mathsf H\)\texttt{TcAnti.lean}, lines 55--71 and 94--100,
and \texttt{Distinct.lean}, lines 21--24.
% END NTL14

% BEGIN NTL15
\paragraph{NTL15. The distinguished coordinate has the required prime and is away from \(\ell\).}
\label{claim:NTL15}
Put \(w_0=w(v_0)\).  Then \(w_0\mid p\), \(w_0\nmid\ell\),
and for every \(u\in U(K)\),
\[
 \operatorname{cmp}_{w_0}(u)\in C_{w_0}(1).
\]
First, \(w(v)\cap\mathcal O_F=v\) for every \(v\).  The two
original fields \(s.\texttt{hw}\ v\) and \(s.\texttt{hA}\ v\)
say that the same valuation subring of \(\Omega\) lies over
\(w(v)\) and \(v\); pulling back its maximal ideal along
\(\mathcal O_F\to\mathcal O_L\to\Omega\) proves the equality.
The residue field at \(v_0\) has cardinality \(p\) by
\eqref{eq:ntl-aux-data}; its characteristic is therefore \(p\),
so \(p\in v_0\), and hence \(p\in w_0\).
If \(\ell\) also belonged to \(w_0\), Bezout's identity for the
distinct rational primes \(p,\ell\) would put \(1\) in \(w_0\).
Thus \(w_0\nmid\ell\).

All of \(U(K)\) is, by construction, a subgroup of \(U_b\).
The second defining condition of \(U_b\) is exactly the displayed
membership in \(C_{w_0}(1)\), independently of \(K\).
Moreover \(s.p\) satisfies the oddness and discriminant hypotheses
of NTL5.  This checks all its local hypotheses at the precise
coordinate used by the original construction.
\emph{Source:} \texttt{nt\_under} and \texttt{nt\_nat}, lines
324--332; auxiliary-place calculation in \texttt{neatU}, lines
434--448; original setting fields in
\texttt{WeilGroupCharacters/}\allowbreak\texttt{ContinuousResP.lean},
lines 73--100.
% END NTL15

% BEGIN NTL16
\paragraph{NTL16. The full stabilizer relation forces its rational matrix to have finite order.}
\label{claim:NTL16}
Assume the antecedent of \eqref{eq:ntl-neat}, and let \(M\) be the
matrix of NTL12.  Then \(M\) has finite order.

Write \(g=(g_a,(g_v)_{v\mid\ell})\) and
\(b=(b_v)_{v\mid\ell}\).  At every \(w\nmid\ell\), projection
of the stabilizer equality and NTL12 give
\[
 M_w\operatorname{cmp}_w(g_a)
   =\operatorname{cmp}_w(g_a)\operatorname{cmp}_w(u).
\]
At a place \(v\mid\ell\), its second coordinate gives
\(M_{w(v)}g_v=g_vb_v\).  Taking powers preserves the conjugation
identities, so for every integer \(k\ge0\),
\begin{equation}\label{eq:ntl-local-conj}
 \begin{aligned}
 \operatorname{cmp}_w(g_a)^{-1}(M^k)_w
       \operatorname{cmp}_w(g_a)&=\operatorname{cmp}_w(u)^k
                              &&(w\nmid\ell),\\
 g_v^{-1}(M^k)_{w(v)}g_v&=b_v^k &&(v\mid\ell).
 \end{aligned}
\end{equation}
By NTL13, \(\operatorname{cmp}_{w_q}(u)\in K_q\).
Apply NTL6 to each of the finitely many \(q\), obtaining
\(e_q>0\) with
\(\operatorname{cmp}_{w_q}(u)^{e_q}\in K_{w_q}^{\max}\).
Put \(N=\prod_{q\in Q}e_q>0\), with empty product \(1\).
Every \(e_q\) divides \(N\).

For every finite place \(v\) of \(F\), we now exhibit a place
\(T(v)\mid v\) at which \(M^N\) is conjugate into the integral
matrix group.  If \(v\mid\ell\), take \(T(v)=w(v)\).  NTL14
and the second equality in \eqref{eq:ntl-local-conj} apply.
For \(v\nmid\ell\), consider the remaining alternatives in
\eqref{eq:ntl-Q}.
Every place above such a \(v\) is also away from \(\ell\), by
contraction of the rational integer \(\ell\).
If \(v\in S_t\), take \(T(v)=w(v)\): membership of \(u\) in
\(U_b\) gives integrality, and the first equality in
\eqref{eq:ntl-local-conj} applies.
If \(v=v_0\), use \(T(v)=w_0\) and NTL15, noting
\(C_{w_0}(1)\subseteq K_{w_0}^{\max}\).
If \(v=v_q\), use \(T(v)=w_q\).  This place is away from
\(\ell\), since \(v_q\notin S\), and the definition of \(N\)
makes \(\operatorname{cmp}_{w_q}(u)^N\) integral.
Finally, if \(v\notin V.S'\), take \(T(v)=w(v)\) and use
\eqref{eq:ntl-level-out}.  These alternatives exhaust all places.
All chosen places really lie over \(v\): for \(w(v)\) this is
NTL15's pullback check; for \(w_q\) it is the definition of \(v_q\).

The matrix \(M^N\) still satisfies the unitarity identity for all
its powers.  NTL11 therefore gives a positive exponent \(a\) with
\((M^N)^a=1\).  Since \(Na>0\), this proves that \(M\) has finite
order.  The finitely many exponents \(e_q\) were chosen for this
stabilizer relation.  No uniform exponent over all \(K\), all \(g\),
or all stabilizers is asserted or needed.
\emph{Source:} the coordinate identities and place-by-place argument
in \texttt{neatU}, lines 417--477.
% END NTL16

% BEGIN NTL17
\paragraph{NTL17. Full original \texttt{neatU}, for every original auxiliary level.}
\label{claim:NTL17}
For every set of source arguments
\[
 \ell,L,\Omega,s,n\ge1,V,Q,c,m,\tau _0,
 \qquad K:(\texttt{hset}\ Q\ \tau _0).\texttt{Lev},
\]
with exactly the structures and hypotheses stated at the beginning
of this chapter, one has
\[
 (\texttt{hset}\ Q\ \tau _0).X.\texttt{Neat}
            \bigl((\texttt{hset}\ Q\ \tau _0).\texttt{lev}\ K\bigr).
\]
In particular this holds for every integral \(\tau _0\), every
\(m\ge0\), and every family of original compact open \(K_q\).

To prove it, take any stabilizer relation
\eqref{eq:ntl-neat}.  By NTL16 its rational matrix \(M\) has finite
order.  NTL15 puts \(\operatorname{cmp}_{w_0}(u)\) in the exact
principal congruence group to which NTL5 applies.  It too has finite
order: the first equality of \eqref{eq:ntl-local-conj}, valid at
\(w_0\nmid\ell\), expresses it as a conjugate of \(M_{w_0}\).
Thus \(\operatorname{cmp}_{w_0}(u)=1\).  The same equality at
exponent one now gives \(M_{w_0}=1\).  The field map
\(L\longrightarrow L_{w_0}\) is injective, so entrywise \(M=1\).
Its rational adelic image is \(\gamma\), whence \(\gamma=1\).
Cancelling \(g\) in the original stabilizer equality gives
\((u,b)=(1,1)\), which is exactly the required conclusion.

The proof has used the fixed auxiliary prime \(s.p\), the actual
groups and maps of \(\texttt{hset}\), and the same input \(K\).
No smoothness or \texttt{TypeForAlt} property of \(\tau _0\), no
\(H.\texttt{Ax}\), and no ordinary module or ideal is presumed.
The source restriction \(n\ge1\) is retained.  The prime hypothesis
\(\ell\ne2\) is already part of \(s\); no inequality between
\(\ell\) and \(n\) has entered.  In particular the permitted
rank-four case \(\ell=3\) is included.
\emph{Source:} \texttt{UpperTheory.V.Hi.neatU}, lines 412--487;
the exact predicate is \(\mathsf H\)\texttt{ContinuousTorusCharacter.lean},
lines 165--166.
% END NTL17

\subsection{What this removes from the Count-II frontier}

NTL17 is an unconditional supplier of the original neatness predicate
from its printed original datum.  Its universal quantifier over
\(K\) is suitable for constructing one torus character before proving
the \texttt{Free} condition at every level.  The character construction
and the elementary implication from \texttt{Neat} to \texttt{Free}
are handled by the companion ordinary-weight chapter; neither has
been smuggled into an assumption here.

No original ideal \(J\) is replaced by this result.  Both rank-four
Links can use NTL17 for their own original data, after the same fixed
\(\iota\) and within the same eventual coefficient range.  There is
no new lower bound on the coefficient prime and no new eventual
threshold.  The existential powers used inside NTL16 depend only
on an individual stabilizer calculation and do not alter any outer
quantifier in the Count-II contract.

The integral type supplier, the remaining original \(H.\texttt{Ax}\)
fields, ordinary modules and control, arithmetic fibers, the four
remaining Hida endpoints, and all-level \texttt{HasAut} on the same
ideal remain separate obligations.  NTL17 closes the full neat-level
producer, not the compound \texttt{COUNT\_II\_SUPPORT} leaf.  Every
claim above is an independently reconstructed elementary or moderate
argument; no published theorem is used as a new class-A import.
